\label{sec:conclusion}

본 논문은 NPU가 없는 초소형 MCU(STM32F746, SRAM 320KB, Flash
1MB)에서 QAT로 학습된 SNN을 정수 전용 코드로 lowering하는
CHACHA 컴파일러와, Conv+AvgPool+LIF를 하나로 묶은 TFLM 커스텀 연산자를
제시하였다. 표준 TFLM+CMSIS-NN 경로가 s16 DSP fast 경로의 커버리지
제약으로 스칼라 구현보다 약 2.1배 느려지는 성능 역전을 실측으로
규명하였고, 제안 퓨전 연산자가 이를 동일 조건 4.0배 가속(이미지당
$2{,}407$ms, 타임스텝당 77.9ms, 아레나 실사용 80.3KB)으로 회복함을
보였다. TFLM 런타임 오버헤드는 약 0.5\%로 무시 가능함을 분리
실측하였으며, 호스트 비트 일치 검증과 LIF SIMD 전수 검증으로 구현의
비트 정확성을 확인하였다. 나아가 메모리 비용 모델과 실측 처리량으로
배포 가능 영역을 정량화하여, 현 보드에서 채널 2배 구성까지 배포
가능(추정)하고 4배 구성은 Flash 초과로 불가함을 제시하였다.

향후 과제는 다음과 같다. (1) 다중 모델$\cdot$다중 데이터셋으로의 검증
확대, (2) 정수 구현의 전체 테스트셋 정확도 실측, (3) 출력 채널 타일
제약($C_o\le64$)의 완화, (4) 타임스텝 축소와 조기 판정에 따른
지연--정확도 트레이드오프 탐색, (5) 전력 실측과 스파이크 희소성을
활용하는 zero-skipping 커널 --- 본 구현은 희소성을 활용하지 않는
dense 연산이며 에너지 실측도 수행하지 않았다 ---, (6)
NIR\cite{pedersen2024nir}와 같은 SNN 교환 포맷으로부터의 임포트
연계이다. 현 컴파일러의 입력은 snnTorch+Brevitas 그래프에 한정되므로,
교환 포맷 연계는 적용 범위를 넓히는 다음 단계가 될 것이다.
